\section{Grundlagen}
\subsection{IFC-Standard}
Der IFC Standard (Industry Foundation Classes) ist eine offene, internationale Spezifikation zur Beschreibung von Bauelementen und deren Beziehungen 
zueinander. Er wird seit 1994 von der Organisation buildingSMART im \gls{Express} Schema entwickelt, um die Digitalisierung auf Baustellen zu 
vereinheitlichen. Er wird von den meisten \ac{BIM} Softwarelösungen verwendet, welche ihre Modellierungen als .ifc Datei exportieren können.
Der IFC Standard besteht aus Objekten, Attributen und Beziehungen. Die Objekte sind hierarchisch aufgebaut. Jedes Objekt 
startet mit der Klasse \textit{IfcRoot} und erbt weitere Attributen von allgemeineren Basisklassen. \
Das Beispiel zeigt die Hierarchie der Klasse \textit{IfcWall} auf:
\begin{figure}[H]
    \centering
    \includegraphics[width=1\textwidth]{IfcWallHierarchie.png} 
    \caption{Hierarchischer Aufbau der IFC-Klassen.}
    \label{fig:wall_hier}
\end{figure}

Die Objekte werden mittels Beziehung in Verbindung gebracht. Zum Beispiel hängt ein Gebäude (\textit{IfcBuilding}) mittels der Relation \textit{IfcRelAggregates}
an einem Bauplatz (\textit{IfcSite}). Als visuelles Beispiel soll folgende Grafik dienen:
\begin{figure}[H]
    \centering
    \includegraphics[width=1\textwidth]{LocationStructrue.png} 
    \caption{Struktueller Aufbau der IFC-Klassen.}
    \label{fig:loc_structure}
\end{figure}

\subsection{Graphdatenbanken}
Graphdatenbanken organisieren ihre Daten, wie der Name schon vermuten lässt, in Graphen. Das bedeutet Objekte werden als Knoten dargestellt, verbunden durch Kanten. Der große Vorteil hierbei ist die native Darstellung von Beziehungen zwischen den Objekten.
Dies wird insbesondere wichtig bei Datenmengen, bei welchen die Beziehungen zueinander ein elementarer Bestandteil der Datenerhebung darstellt. Bei den gängigsten Graphdatenbanken handelt es sich um sogenannte Labled Property Graphs. Hierbei können nicht nur
Knoten, sondern auch Kanten mit Labels und Eigenschaften versehen werden. 
\begin{figure}[htbp]
    \centering
    \includegraphics[width=0.5\textwidth]{property-graph.png} 
    \caption{Labled Property Graph \newline 
    (\url{https://dist.neo4j.com/wp-content/uploads/20240604025137/property-graph-model-1.png})}
    \label{fig:property_graph}
\end{figure}
Dieser Ansatz bietet eine hohe Flexibilität gegenüber relationalen Datenbanken, da hier der Datensatz logisch miteinander verknüpft werden kann und aufwände JOIN Operationen vermieden
werden können \cite{seidel2020graph}. Graphdatenbanken erfreuen sich vor allem großer Beliebtheit bei stark vernetzten Daten, wie Social Media, Betrugerkennung oder bei Lieferketten \cite{google_graphdb}, da hierbei der Analyse der Verbindungen im Vordergrund
steht.
%%TODO Eventuell Umschreiben (letzter Satz)
\subsubsection{neo4j}
Neo4j ist eine \gls{acid} konforme Graphdatenbank auf der Grundlage von Labled Property Graphs. Sie wurde erstmalig 2007 von Emil Eifrem veröffentlicht und avancierte zu einer der weltweit führenden Graphdatenbanken.

\subsubsection{GraphQL}
\subsection{Mircofrontends}
\subsection{Monorepos}